% Problem record op_eba8dda5ff69d0f2. This TeX file is derived from
% database/problems_json/op_eba8dda5ff69d0f2.json by scripts/sync-tex.mjs; edit the JSON record
% and rerun the script rather than editing this file.
\section{Minimum output entropy additivity for qubit channels}
\label{sec:op_eba8dda5ff69d0f2}

\paragraph{Problem.}
Is minimum output von Neumann entropy additive whenever one channel is a qubit channel? Let $\Phi:\mathcal L(\mathbb C^2)\to\mathcal L(\mathbb C^2)$ be any completely positive trace-preserving map, and let $\Omega:\mathcal L(A)\to\mathcal L(B)$ be any finite-dimensional quantum channel. For a channel $\mathcal N$, define
\begin{equation}
 S_{\min}(\mathcal N):=\min_{\rho\in\mathcal D(A_{\mathcal N})}S(\mathcal N(\rho)),\qquad S(\sigma):=-\operatorname{Tr}\sigma\log_2\sigma,
 \label{eq:h11-quantities}
\end{equation}
where $\mathcal D(A_{\mathcal N})$ is the set of density operators on the input space of $\mathcal N$. Prove or disprove
\begin{equation}
 S_{\min}(\Phi\otimes\Omega)=S_{\min}(\Phi)+S_{\min}(\Omega)\qquad\text{for every such pair }(\Phi,\Omega).
 \label{eq:h11-additivity}
\end{equation}
In Eq.~\eqref{eq:h11-additivity}, both the input and output of $\Phi$ have dimension two; no unitality assumption $\Phi(I_2)=I_2$ is imposed. Since product inputs are admissible in the minimization of Eq.~\eqref{eq:h11-quantities}, a counterexample must certify strict subadditivity for a specific finite-dimensional pair.

\subsection*{Status}

Unsolved

\subsection*{Source}

Holevo explicitly records the unresolved nonunital-qubit case in Section~2.5, p.~13 of the arXiv version, in a discussion where the second channel is arbitrary \sourcecite{ref:h11-holevo}{Hol15}. Hayashi reviews the qubit numerical evidence and the unital theorem in Section~9.9.1, pp.~552--553 \sourcecite{ref:h11-book}{Hay17}. Ruskai notes the unresolved nonunital self-power possibility immediately after Problem~19, p.~16 \sourcecite{ref:h11-ruskai}{Rus07}.

\subsection*{Progress}

\begin{itemize}
  \item King's Theorem~1 proves the equality in Eq.~\eqref{eq:h11-additivity} for every unital qubit channel with an arbitrary channel as partner. It also proves maximal output Schatten-$p$ norm multiplicativity for every $p\geq1$ in this unital setting \sourcecite{ref:h11-king}{Kin02}.

  \item King and Koldan's Theorem~1 proves maximal output Schatten-$p$ norm multiplicativity for arbitrary qubit channels at $p=2$ and $p\geq4$, with an arbitrary finite-dimensional completely positive partner. These orders do not establish the von Neumann-entropy equality in Eq.~\eqref{eq:h11-additivity} \sourcecite{ref:h11-kk}{KK06}.

  \item Large-dimensional nonadditivity results do not settle the qubit restriction. Leung, Lovitz, and Wu's July~2026 preprint obtains random-channel counterexamples for R\'enyi orders $p>3/4$ and $0\leq p<1/4$. At $p=1$, its Section~5.1 numerically evaluates an ensemble-specific Bell-input threshold at output dimension $182$; this is neither a qubit counterexample nor a proof of a universal dimension threshold \sourcecite{ref:h11-llw}{LLW26}.
\end{itemize}

\subsection*{References}

\begin{enumerate}[label={},leftmargin=4.8em,labelsep=0.5em]
  \footnotesize
  \item[\textup{[Hay17]}]\label{ref:h11-book}
  M. Hayashi, \emph{Quantum Information Theory: Mathematical Foundation}, 2nd ed., Graduate Texts in Physics, Springer (2017). \href{https://doi.org/10.1007/978-3-662-49725-8}{doi:10.1007/978-3-662-49725-8}.

  \item[\textup{[Hol15]}]\label{ref:h11-holevo}
  A. S. Holevo, ``Gaussian optimizers and the additivity problem in quantum information theory,'' \emph{Russian Mathematical Surveys} \textbf{70}(2), 331--367 (2015). \href{https://doi.org/10.1070/RM2015v070n02ABEH004949}{doi:10.1070/RM2015v070n02ABEH004949}; \href{https://arxiv.org/abs/1501.00652}{arXiv:1501.00652}.

  \item[\textup{[Rus07]}]\label{ref:h11-ruskai}
  M. B. Ruskai, ``Open Problems in Quantum Information Theory,'' preprint (2007). \href{https://arxiv.org/abs/0708.1902}{arXiv:0708.1902}.

  \item[\textup{[Kin02]}]\label{ref:h11-king}
  C. King, ``Additivity for unital qubit channels,'' \emph{Journal of Mathematical Physics} \textbf{43}, 4641--4653 (2002). \href{https://doi.org/10.1063/1.1500791}{doi:10.1063/1.1500791}; \href{https://arxiv.org/abs/quant-ph/0103156}{arXiv:quant-ph/0103156}.

  \item[\textup{[KK06]}]\label{ref:h11-kk}
  C. King and N. Koldan, ``New multiplicativity results for qubit maps,'' \emph{Journal of Mathematical Physics} \textbf{47}(4), 042106 (2006). \href{https://doi.org/10.1063/1.2191787}{doi:10.1063/1.2191787}; \href{https://arxiv.org/abs/quant-ph/0512185}{arXiv:quant-ph/0512185}.

  \item[\textup{[LLW26]}]\label{ref:h11-llw}
  D. Leung, B. Lovitz, and P. Wu, ``Counterexamples to additivity of minimum output $p$-R\'enyi entropy of quantum channels for all $p>3/4$ and $0\leq p<1/4$,'' preprint (2026). \href{https://arxiv.org/abs/2607.15210}{arXiv:2607.15210}.
\end{enumerate}

\subsection*{Comment}

Audited on 2026-09-09 against the full source statements and later nonadditivity literature. The nonunital qubit case remains unresolved. The smallest-output-dimension record permits arbitrary input dimensions and bounds both output dimensions, whereas this question fixes one channel to have qubit input and output and leaves its partner unrestricted. The closed-form-witness record permits arbitrary dimensions, so a solution there need not answer this question. Two related unresolved variants are Holevo-capacity additivity with an arbitrary partner and minimum-output-entropy additivity on self-powers of a qubit channel. A general partner counterexample need not yield a qubit self-power counterexample. The universal additivity equivalences in unrestricted dimensions do not justify identifying the Holevo and minimum-output-entropy questions with these dimension constraints.

\subsection*{Field}

Quantum Communication

\subsection*{Topic}

Additivity and regularization; Matrix and entropy inequalities; Quantum channel structure

\subsection*{ID}

\texttt{op\_eba8dda5ff69d0f2}
